\section{Conclusiones}
El desarrollo del sistema de gestión de tickets confirma que la integración de frameworks modernos y metodologías ágiles constituye una estrategia efectiva para optimizar procesos de atención de incidencias. La elección de Laravel y React permitió construir un sistema robusto, escalable y centrado en la experiencia de usuario, mientras que Scrum facilitó la adaptabilidad del proceso de desarrollo.

Entre los principales aportes se destacan: la centralización de información, la reducción en los tiempos de atención, la personalización de roles y la generación de métricas confiables para la toma de decisiones. Asimismo, el sistema sienta las bases para futuras extensiones, incluyendo compatibilidad multiplataforma con React Native, integración de microservicios y el uso de inteligencia artificial en la clasificación de tickets.

En conclusión, el proyecto valida la pertinencia de unir teoría y práctica en el desarrollo de soluciones tecnológicas, demostrando que la literatura científica puede guiar de manera efectiva la toma de decisiones en proyectos reales.